\documentclass[11pt]{article}

\usepackage{amsmath, amssymb, amsthm, geometry}
\geometry{margin=1in}

\title{A Layered Framework for Spectral, Geometric, and Symbolic Representations of Zeta and Prime Structures}
\author{Redacted for Review}
\date{}

% ---------------- Theorem Environments ----------------
\newtheorem{definition}{Definition}
\newtheorem{theorem}{Theorem}
\newtheorem{proposition}{Proposition}
\newtheorem{remark}{Remark}
\newtheorem{axiom}{Axiom}

\begin{document}

\maketitle

% ======================================================
\section{Abstract}

This paper introduces a three-layer framework for analyzing mathematical structures associated with the Riemann zeta function, prime distributions, and geometric embeddings. The framework explicitly separates formal mathematical objects, computational visualization mappings, and symbolic interpretation systems. This separation is introduced to avoid conflating proven mathematical structures with metaphorical or representational constructs.

% ======================================================
\section{Axiomatic Foundation}

\begin{axiom}[Epistemic Separation Principle (ESP)]
No symbolic, geometric, or computational representation shall be treated as mathematically equivalent to a formal object unless it is derived through explicit axioms, definitions, and proofs within a shared mathematical structure.
\end{axiom}

\begin{remark}
This axiom ensures that interpretive or visual representations cannot be mistaken for ontological or analytic equivalence.
\end{remark}

% ======================================================
\section{Layer I: Formal Mathematical Structures}

\subsection{Zeta Function and Zeros}

\begin{definition}[Riemann Zeta Function]
For $\mathrm{Re}(s) > 1$, the Riemann zeta function is defined as:
\[
\zeta(s) = \sum_{n=1}^{\infty} \frac{1}{n^s}
\]
with meromorphic continuation to $\mathbb{C} \setminus \{1\}$.
\end{definition}

\begin{definition}[Nontrivial Zeros]
Nontrivial zeros of $\zeta(s)$ are complex numbers $\rho_n = \beta_n + i t_n$ satisfying:
\[
\zeta(\rho_n) = 0, \quad 0 < \beta_n < 1.
\]
\end{definition}

\begin{conjecture}[Riemann Hypothesis]
For all nontrivial zeros $\rho_n$ of $\zeta(s)$:
\[
\beta_n = \frac{1}{2}.
\]
\end{conjecture}

% ======================================================
\subsection{Spectral Operator Hypothesis}

\begin{conjecture}[Hilbert--P\'{o}lya Type Operator]
There exists a self-adjoint operator $H$ on a Hilbert space $\mathcal{H}$ such that:
\[
H \psi_n = t_n \psi_n,
\]
where $t_n$ are the imaginary parts of the nontrivial zeros of $\zeta(s)$.
\end{conjecture}

\begin{remark}
This statement is unproven and treated as a structural hypothesis in spectral number theory.
\end{remark}

% ======================================================
\subsection{Graph-Theoretic Structure}

\begin{definition}[Weighted Graph Model]
Let $G = (V, E, W)$ be a weighted graph with adjacency matrix $W_{ij}$ and degree matrix $D$. The graph Laplacian is defined as:
\[
L = D - W.
\]
\end{definition}

\begin{proposition}
The Laplacian $L$ is symmetric and positive semi-definite if $W$ is symmetric and non-negative.
\end{proposition}

% ======================================================
\section{Layer II: Computational Embedding Model}

\begin{definition}[Geometric Embedding Map]
Let $\mathcal{F}$ be a mapping:
\[
\mathcal{F}: \mathcal{D} \rightarrow \mathbb{R}^2
\]
from discrete mathematical objects $\mathcal{D}$ into Euclidean coordinates for visualization purposes.
\end{definition}

\begin{remark}
This mapping does not preserve algebraic or spectral structure; it is purely representational.
\end{remark}

\subsection{Logarithmic Spiral Embedding}

A common embedding used in this framework is:
\[
r(\theta) = \phi^{\theta},
\]
where $\phi = \frac{1+\sqrt{5}}{2}$ is the golden ratio.

\begin{remark}
This function is used only for visualization scaling and does not encode physical or spectral equivalence.
\end{remark}

% ======================================================
\section{Layer III: Symbolic Interpretation System}

\begin{definition}[Symbolic Mapping Layer]
A symbolic mapping is a function:
\[
\mathcal{S}: \mathcal{D} \rightarrow \Sigma,
\]
where $\Sigma$ is a set of interpretive labels (e.g., frequencies, archetypes, or metaphoric constants).
\end{definition}

\begin{remark}
Symbolic mappings are explicitly non-mathematical in the sense that they do not define or constrain formal structures in Layer I.
\end{remark}

% ======================================================
\section{Separation Theorem}

\begin{theorem}[Layer Independence]
Let $\mathcal{L}_1, \mathcal{L}_2, \mathcal{L}_3$ denote Layers I--III. Then:
\[
\mathcal{L}_1 \not\equiv \mathcal{L}_2, \quad
\mathcal{L}_2 \not\equiv \mathcal{L}_3, \quad
\mathcal{L}_1 \not\equiv \mathcal{L}_3.
\]
\end{theorem}

\begin{proof}
Each layer operates on distinct ontological assumptions:
\begin{itemize}
\item $\mathcal{L}_1$: formal axiomatic mathematics
\item $\mathcal{L}_2$: coordinate embeddings without structural preservation guarantees
\item $\mathcal{L}_3$: interpretive symbolic mapping with no constraint preservation
\end{itemize}
Therefore, no isomorphism is defined or implied between layers.
\end{proof}

% ======================================================
\section{Discussion}

This framework allows exploration of:
\begin{itemize}
\item spectral properties of zeta zeros
\item geometric visualization of number-theoretic objects
\item symbolic interpretation systems for pattern generation
\end{itemize}

However, only Layer I supports formal mathematical claims.

% ======================================================
\section{Limitations}

\begin{itemize}
\item No claim is made that symbolic or geometric representations reveal hidden physical laws.
\item No equivalence between frequency-based interpretations and mathematical constants is assumed.
\item The framework is not predictive in physical or cosmological domains.
\end{itemize}

% ======================================================
\section{Conclusion}

We propose a disciplined separation of mathematical structure, computational visualization, and symbolic interpretation. This separation allows exploratory modeling while preserving mathematical rigor and preventing category errors between representation and formal truth.

\end{document}
